\graphicspath{{abstracts/praesentationen/04-sparql-endpunkte/img/}{img/}}

\begin{beitrag}{SPARQL-Endpunkte als nachhaltige Schnittstelle für geisteswissenschaftliche Forschungsdaten}{%
Greta Gärtner\,\orcidlink{0000-0001-8901-2345} \\
\small \href{mailto:greta.gaertner@example.org}{greta.gaertner@example.org} \\
\small Friedrich-Alexander-Universität Erlangen-Nürnberg \\[0.5cm]
Hannes Hauser\,\orcidlink{0000-0002-9012-3456} \\
\small \href{mailto:hannes.hauser@example.org}{hannes.hauser@example.org} \\
\small Technische Universität Dresden
}

\begin{abstract}
Der Beitrag berichtet über Aufbau und Betrieb stabiler, öffentlich zugänglicher SPARQL-Endpunkte für vier WissKI-Projekte über drei Jahre: Die Architektur trennt die redaktionelle WissKI-Instanz strikt vom Veröffentlichungsindex (NGINX, OpenLink Virtuoso, Apache Airflow), der automatisch mit VoID-Beschreibungen und DataCite-konformen Metadaten versehen wird. Das Service-Modell kommuniziert Verfügbarkeit als "best effort with disclosure" und überwacht kritische Schwellenwerte mit Prometheus. Nach drei Betriebsjahren zeigt sich, dass SPARQL-Endpunkte ebenso als didaktisches Werkzeug wie als Forschungsschnittstelle genutzt werden und der Aufwand für den Eigenbetrieb geringer ist als oft angenommen.
\end{abstract}

\section{Motivation}
Mit dem zunehmenden Einsatz semantischer Datenmodelle in der geisteswissenschaftlichen Forschung wird die Frage der Nachnutzbarkeit dringlicher. Datenexporte als CSV oder XML können zwar vollständig sein, sie sind aber statisch, verlieren ihre semantische Tiefe und sind nach wenigen Jahren häufig nicht mehr aktuell. Ein vom Projekt selbst betriebener SPARQL-Endpunkt verspricht eine Alternative: Forschungsdaten bleiben in ihrer ursprünglichen Modellierung, sind über Standardschnittstellen abfragbar und können auch von Dritten in eigenen Forschungskontexten verknüpft werden. Vor diesem Hintergrund haben wir im Verlauf der vergangenen drei Jahre für vier WissKI-Projekte stabile, öffentlich zugängliche SPARQL-Endpunkte aufgebaut und in einer institutionellen Infrastruktur konsolidiert.

\section{Technische Architektur}
Die Endpunkte werden über eine schlanke Frontend-Schicht (NGINX, OpenLink Virtuoso) ausgespielt, hinter der die einzelnen Projekt-Triplestores in benannten Graphen zusammengefasst sind. Eine zentrale Designentscheidung war die strikte Trennung zwischen redaktioneller WissKI-Instanz und dem öffentlich verfügbaren Endpunkt: Inhalte werden nicht direkt aus dem Drupal-Frontend bedient, sondern in regelmäßigen Abständen aus dem Bearbeitungssystem in den Veröffentlichungsindex repliziert. Diese Zweischichtigkeit hat sich als wesentliche Voraussetzung sowohl für die Performance unter Last als auch für die Steuerung der Sichtbarkeit erwiesen; ausführlich diskutiert wird der Architekturentscheid bei \cite{gaertner2024-sparql-architektur}.

Die Replikation erfolgt über einen mit Apache Airflow orchestrierten Workflow, der zusätzlich projektabhängige Aufbereitungsschritte ausführt: Berechnung abgeleiteter Klassifizierungen, Anreicherung um Normdaten und das Schreiben einer kompakten VoID-Beschreibung des publizierten Datensatzes. Diese VoID-Datei wird zusammen mit einer maschinenlesbaren Lizenz und einem DataCite-konformen Metadatensatz veröffentlicht, sodass die Endpunkte über Standardregister auffindbar sind.

\section{Service-Level und Monitoring}
Ein häufiger Einwand gegen den Eigenbetrieb von SPARQL-Endpunkten lautet, dass die Forschungsteams nicht in der Lage seien, einen produktiven Dienst langfristig zuverlässig zu betreiben. Wir haben deshalb von Beginn an auf ein minimal-invasives, aber explizit dokumentiertes Service-Modell gesetzt. Verfügbarkeit, Antwortzeit und Anfragezahl werden mit Prometheus überwacht, kritische Schwellenwerte führen zu Benachrichtigungen an die IT-Servicegruppe der Universität. Die Endpunkte werden ausdrücklich nicht mit einer harten Verfügbarkeitsgarantie versehen, sondern als \glqq best effort with disclosure\grqq{} kommuniziert -- ein realistischer Anspruch, der den tatsächlichen Möglichkeiten der akademischen Infrastruktur entspricht.

Auf Anwendungsseite stellen die Endpunkte sowohl eine klassische SPARQL-Schnittstelle als auch eine kuratierte Sammlung vorgefertigter Abfragen zur Verfügung. Diese vorgefertigten Abfragen sind nicht nur ein Werkzeug für die Wissensvermittlung, sondern auch ein Test-Set für die kontinuierliche Validierung: Wenn sich nach einer Aktualisierung das Ergebnis einer dieser Abfragen ändert, wird das automatisch protokolliert und vom Projektteam geprüft. Eine systematische Auswertung der so erfassten Veränderungen findet sich bei \cite{hauser2025-endpunkt-monitoring}.

\section{Erfahrungen}
Nach drei Betriebsjahren ergibt sich folgendes Bild: Die Endpunkte werden überwiegend von der eigenen Projektgemeinschaft und einer kleinen Zahl externer Forschender genutzt; massive Lastspitzen sind die seltene Ausnahme. Die wichtigste Form der Nachnutzung liegt im niedrigschwelligen Zugriff auf Lehre und Demonstration -- der Endpunkt ist ein didaktisches Werkzeug ebenso wie ein Forschungsdienst. Wir empfehlen daher, in der Konzeption solcher Endpunkte beide Nutzungsformen explizit zu berücksichtigen.

\section{Fazit}
SPARQL-Endpunkte sind ein lohnender Baustein der eigenen Forschungsdateninfrastruktur, wenn sie als Teil einer klar dokumentierten Service-Architektur betrachtet werden. Der Aufwand ist überschaubar, der methodische Mehrwert hoch, und die Hürden für die institutionelle Verstetigung sind geringer, als oftmals angenommen wird.

\end{beitrag}

\section*{Kurzbiografie(n)}

\subsection*{Greta Gärtner}
Greta Gärtner ist Postdoktorandin an der FAU Erlangen-Nürnberg und Mitglied des WissKI-Konsortiums. Ihre Forschungsschwerpunkte liegen in der Anwendung semantischer Datenmodelle auf museale Sammlungen mit kleinem und mittelgroßem Bestand. Sie hat in mehreren Migrationsprojekten als wissenschaftliche Begleitung mitgewirkt und beschäftigt sich aktuell mit der Schnittstellenanbindung von WissKI-Beständen an Museums-Verbund-Plattformen.

\subsection*{Hannes Hauser}
Hannes Hauser ist Forschungsdatenmanager am Sächsischen Landesdigitalisierungszentrum an der TU Dresden. Er betreut SPARQL-Endpunkte und Datenpublikations-Workflows für mehrere geisteswissenschaftliche Forschungsprojekte und entwickelt im Rahmen einer Promotionsarbeit ein Monitoring-Framework für die Langzeitstabilität semantischer Forschungsdatendienste.
